\documentclass[a4paper]{article}
\usepackage[utf8]{inputenc}
\usepackage[english]{babel}
\usepackage{datetime}
\usepackage{hyperref}
\usepackage[a4paper, margin=2cm]{geometry}

\newdateformat{monthyeardate}{\monthname[\THEMONTH] of \THEYEAR}

\pagestyle{empty}
\setlength{\parindent}{0pt}

\begin{document}

\begin{center}
  {\Large \textbf{Francisco Figari}} \\[4pt]
  Buenos Aires based Software Developer \\[2pt]
  {\small \href{mailto:figafran@gmail.com}{figafran@gmail.com} -- \href{https://www.linkedin.com/in/ffigari}{linkedin.com/in/ffigari}} \\[2pt]
  {\small \textit{Resume} -- \monthyeardate\today}
\end{center}

\vspace{0.5em}

\section*{Professional Summary}

% The TPO summary emphasizes bridging product and tech. For a developer resume,
% the summary should foreground your hands-on building experience.
% Questions to answer here:
% - What kind of systems have you built? (microservices, web apps, simulations, etc.)
% - What is your strongest technical identity? (backend dev? systems thinker? polyglot?)
% - What sets you apart from other developers with similar years of experience?
% - How does your cross-domain experience (fintech, urbanism, neuroscience) make
%   you a better developer specifically? (e.g., tackling novel problem domains)
With ~9 years of experience in the software industry, I continuously deliver
implementations accross different domains and teams of different structure.

\section*{Relevant Experience}

\textbf{Technical Leader @ Urbanly} \hfill August 2025 -- present
\begin{itemize}
% In the TPO resume this role is product-focused. For a developer resume,
% highlight what you actually built and the technical decisions you made.
% - What tech stack did you use for the housing simulation PoC?
% - What was the architecture? (monolith, modular, etc.)
% - Were there interesting algorithmic challenges in the simulation?
% - Did you set up CI/CD, testing infrastructure, or dev tooling for the team?
% - What development practices did you introduce and how did they concretely
%   improve the codebase? (e.g., reduced bugs, faster onboarding)
% - Any performance work? Data modeling decisions?
  \item ...
\end{itemize}

\textbf{Backend Developer @ Brubank} \hfill April 2023 -- August 2025
\begin{itemize}
% This is probably your strongest section for a developer resume. The TPO version
% mentions microservices and features but stays high-level. Go deeper here.
%
% Microservices ecosystem:
% - How many services did you personally own or contribute to?
% - What was the tech stack? (Golang? Node/TS? What message broker? What DB?)
% - Did you extract services from the monolith? What was the migration strategy?
% - What scale did your services handle? (the old resume mentions 200 rps --
%   expand on this: peak loads, latency targets, etc.)
%
% Feature development (CEDEARs, stocks, time deposits):
% - What was technically challenging about these features?
% - Did you integrate with external APIs (e.g., market data providers, brokers)?
%   What were the reliability/consistency challenges?
% - How did you handle things like idempotency, retries, eventual consistency?
%
% System quality:
% - How did you use Datadog? Did you set up dashboards, alerts, SLOs?
% - Did you do any performance optimization work?
% - What was your testing strategy? Integration tests? Contract tests?
%
% Interface design:
% - You mention defining mobile-to-backend interfaces. Can you describe a
%   specific API you designed and why the design was good?
%
% Other:
% - Did you do code reviews? What did you focus on?
% - Any incident response or debugging war stories?
% - Did you use feature flags for gradual rollouts? How was that implemented?
  \item ...
\end{itemize}

\textbf{Research Intern @ LIAA, DC, FCEyN, UBA} \hfill August 2021 -- September 2022
\begin{itemize}
% The TPO version is fine here but doesn't mention what you built.
% - What was the tech stack for the web eye tracker? (JS? WebGL? What browser APIs?)
% - What were the main technical challenges? (real-time performance in browser?
%   calibration algorithms? signal processing?)
% - Did you benchmark your implementation against existing tools? What were the results?
% - Was there any data pipeline work (collecting, storing, analyzing eye tracking data)?
  \item ...
\end{itemize}

\textbf{Lead Programmer @ UrbanSim} \hfill July 2018 -- November 2020
\begin{itemize}
% The procedural generator is a strong developer story. Expand on it.
% - What language/framework? (C++? Python? Web-based?)
% - What was the algorithm behind the procedural generation?
% - How did you encode urban planning constraints formally?
% - What was the web platform's stack?
% - Any interesting geometry/computational geometry work?
% - How did you handle the complexity of the constraint system as it grew?
  \item ...
\end{itemize}

% Should any of these older roles appear?
% - DubbingDigital (fullstack): What did you build? What stack?
% - Carleton University (research intern): What did you work on technically?
% - Eryx (junior programmer): Relevant if it shows breadth, but maybe too old.
% - Teaching assistant: Could be left for Education, or dropped.
% Think about which ones add signal for a developer role vs. just adding length.

\section*{Key Competencies}

\subsection*{Software Architecture \& Design}
\begin{itemize}
% This replaces "Product & Delivery". Focus on how you think about building software.
% - How do you approach breaking down a system into components/services?
% - How do you think about API design? What principles guide you?
% - How do you handle growing codebases? (refactoring strategies, decomposition)
% - How do you think about data modeling and persistence choices?
% - Do you have experience with event-driven architectures, message queues, etc.?
  \item ...
\end{itemize}

\subsection*{Technical Stack}
\begin{itemize}
% Be concrete here. What can you actually work with?
% - Which languages are you strongest in? Rank them by proficiency/recency.
% - Frameworks: any specific ones worth calling out? (e.g., Express, Gin, FastAPI)
% - Databases: PostgreSQL, Redis -- at what depth? Schema design? Query optimization?
% - Infrastructure: Docker? Kubernetes? AWS/GCP services? CI/CD tools?
% - Monitoring: Datadog -- what specifically? APM, logs, custom metrics?
% - Testing: what kinds of tests do you write? What tools?
% - AI tooling: Claude Code, copilot -- how do you use them concretely?
  \item ...
\end{itemize}

\subsection*{Collaboration \& Communication}
\begin{itemize}
% Keep this shorter than the TPO version. Focus on what matters for a dev role.
% - Code review culture: what do you bring to it?
% - How do you work with other devs on shared systems?
% - Cross-team work: how do you collaborate with mobile, QA, infra?
% - Mentoring: how have you helped junior devs grow technically?
% - Interdisciplinary background: frame it as "can understand novel domains quickly"
% - Languages: Spanish, French, English
  \item ...
\end{itemize}

\section*{Education and Training}

\begin{itemize}
  \item \textit{Licenciatura} in Computer Science @ \textit{DC, FCEyN, UBA}, equivalent to Bachelor + Master with thesis, started in 2013 and finished in 2022.
  \item High school @ \textit{Lyc\'{e}e Jean Mermoz, Buenos Aires}, finished in 2013.
\end{itemize}

% Should the Product Discovery course stay? It's more product-oriented.
% The yoga training -- keep it? It shows something about you as a person,
% but consider whether it helps or distracts in a developer resume.

\end{document}
